\documentclass{article}
\usepackage[utf8]{inputenc}
\usepackage{geometry}
\geometry{a4paper, margin=1in}
\usepackage{minted}

\title{UAV Software Engineering: ROS 2 Architecture}
\author{Workspaces, Colcon, and Ament}
\date{}

\begin{document}

\maketitle

\section{The Concept of a Workspace}
In standard C++ development, a project typically resides in a single directory with one root \texttt{CMakeLists.txt}. 
In robotics, a UAV software stack is composed of dozens of independent packages (e.g., \texttt{camera\_driver}, \texttt{flight\_controller}, 
\texttt{telemetry\_node}). 

ROS 2 organizes these packages into a \textbf{Workspace}. A standard workspace contains three primary directories:
\begin{itemize}
    \item \texttt{src/}: The only directory you will manually edit. It contains the source code for all your individual packages.
    \item \texttt{build/}: Where intermediate files (object files, CMake cache) are stored for each package.
    \item \texttt{install/}: Where the final compiled executables, libraries, and configuration scripts are placed.
\end{itemize}

\section{Colcon: The Meta-Build System}
To compile a workspace, ROS 2 uses \textbf{Colcon} (COLlective CONstruction). Instead of manually running \texttt{cmake} and \texttt{make} in each package, 
you run \texttt{colcon build} at the root of the workspace. Colcon analyzes the dependency graph of all packages in the \texttt{src/} folder, 
determines the correct build order, and invokes CMake for each one concurrently.

\section{Package Anatomy: \texttt{package.xml} and \texttt{CMakeLists.txt}}
For Colcon to recognize a folder inside \texttt{src/} as a ROS 2 package, it must contain two files:
\begin{enumerate}
    \item \textbf{\texttt{package.xml}:} An XML manifest defining the package's name, version, license, and its dependencies (e.g., \texttt{rclcpp}, 
    \texttt{std\_msgs}). This is crucial for resolving dependencies before the build process begins.
    \item \textbf{\texttt{CMakeLists.txt}:} The build script. However, ROS 2 extends standard CMake with the \textbf{\texttt{ament\_cmake}} macro ecosystem. 
    Ament simplifies finding ROS 2 dependencies and installing the generated executables into the workspace's \texttt{install/} directory.
\end{enumerate}

\section{The Overlay Environment (\texttt{setup.bash})}
Linux needs to know where to find your newly compiled UAV executables and libraries. 
When Colcon finishes building, it generates a \texttt{setup.bash} file inside the \texttt{install/} directory. 

By executing \texttt{source install/setup.bash}, you prepend your workspace's paths to your system's environment variables (like \texttt{PATH} and 
\texttt{LD\_LIBRARY\_PATH}). This process creates an \textbf{Overlay} on top of the base ROS 2 installation (the \textit{Underlay}), allowing the terminal 
to dynamically discover and run your custom drone nodes.

\section{Advanced Colcon Strategies: Symlinks and Build Types}
Building a massive UAV monorepo can take minutes or even hours. To optimize the development cycle, engineers rely on specific Colcon arguments:
\begin{itemize}
    \item \textbf{\texttt{--symlink-install}:} By default, Colcon copies files from \texttt{src/} to \texttt{install/}. 
    If you modify a Python script, a launch file, or a YAML configuration, you are forced to rebuild. Using \texttt{--symlink-install} creates symbolic links 
    instead of copying. This allows changes to interpreted scripts and config files to take effect immediately without recompiling the workspace.
    \item \textbf{CMake Build Types:} By default, CMake compiles without heavy optimizations. In drone software, running unoptimized control loops can cause 
    the CPU to miss real-time deadlines, leading to catastrophic crashes. It is mandatory to compile flight-ready nodes with hardware-specific 
    optimizations using: \\
    \texttt{colcon build --cmake-args -DCMAKE\_BUILD\_TYPE=Release}.
\end{itemize}

\section{The Power of \texttt{ament\_cmake} and Transitive Dependencies}
Standard CMake requires developers to manually track and link every library and its dependencies 
using \texttt{target\_link\_libraries} and \texttt{target\_include\_directories}. In a complex robotics stack, this becomes unmaintainable.

ROS 2 solves this with \textbf{\texttt{ament\_cmake}}. When you use \texttt{ament\_target\_dependencies(my\_node rclcpp std\_msgs)}, 
Ament automatically resolves all \textit{transitive dependencies} (the dependencies of your dependencies). 
Furthermore, when you export your package using \texttt{ament\_export\_dependencies()}, any other node in the UAV stack that depends on your package will 
automatically inherit the correct include paths and linked libraries, ensuring a modular and robust build graph.

\section{Environment Chaining: Underlays and Overlays}
Professional drone software stacks are often containerized (as we will see in Chapter 7). 
Understanding how workspaces chain together is vital. 
The system installation of ROS 2 (e.g., in \texttt{/opt/ros/humble}) is the \textbf{Underlay}. 
When you source its \texttt{setup.bash}, your environment is populated with the base ROS 2 libraries.

When you compile your custom \texttt{uav\_ws} and source its \texttt{setup.bash}, it becomes an \textbf{Overlay}. 
The overlay intelligently prepends its paths to the environment without overwriting the underlay. 
This allows you to have multiple layers of workspaces (e.g., a base ROS 2 underlay, a third-party driver middle-layer, and your custom flight controller overlay),
enabling seamless integration of proprietary and open-source components.

\section{Practical Workflow: Workspace Initialization}
The following terminal commands represent the standard workflow to create a workspace, generate a package, and compile the overlay:

\begin{minted}{bash}
# 1. Create the workspace and the source directory
mkdir -p uav_ws/src
cd uav_ws/src

# 2. Source the system Underlay to enable ROS 2 CLI tools
# (Replace 'jazzy' with your specific ROS 2 distribution)
source /opt/ros/jazzy/setup.bash

# 3. Create a C++ package using ament_cmake
ros2 pkg create --build-type ament_cmake dummy_uav_pkg

# 4. Return to the workspace root and build using Colcon
cd ..
colcon build

# 5. Source the newly built Overlay to make the package visible to Linux
source install/setup.bash
\end{minted}

\section{Advanced Build Optimization}
In production UAV environments, building can be optimized using:
\begin{itemize}
    \item \texttt{--symlink-install}: Creates symbolic links instead of copying files, allowing Python scripts and config files to update immediately without recompiling.
    \item \texttt{Release Mode}: Drone software must meet hard real-time deadlines. Compile with hardware optimizations using: \\
    \texttt{colcon build --cmake-args -DCMAKE\_BUILD\_TYPE=Release}.
\end{itemize}

\end{document}